\subsection{带赋值的域上的拓扑向量空间}

通篇设 K 是一个（未必交换的）域，v 是 K 上的一个赋值，G 是它的阶群。K 赋以拓扑 \(T_{v}\) 。

命题 3. 设 E 是 K 上的一个 Hausdorff 的、维数为 1 的左拓扑向量空间。假设 v 不是非真的。对 E 中一切 \(x_{0} \neq 0\) ，映射 \(a \mapsto ax_{0}\) 把 \(K_{s}\) 映到 E 上，是一个拓扑同构。

这个映射是连续的代数同构。只需证明它是双连续的。设 \(\alpha \in \mathbf{G}\) 。需要证明存在 E 中 0 的一个邻域 V，使得关系 \(ax_0 \in \mathbf{V}\) 蕴含 \(v(a) > \alpha\) 。设 \(a_0 \in \mathbf{K}^*\) 满足 \(v(a_0) = \alpha\) 。由于 E 是 Hausdorff 的，存在 E 中 0 的邻域 W 使得 \(a_0 x_0 \notin \mathbf{W}\) 。由于 \(v\) 不是非真的，存在 E 中 0 的邻域 W' 与 G 的元素 \(\beta\) ，使得关系 \(y \in \mathbf{W}'\) 与 \(v(a) \geqslant \beta\) 蕴含 \(ay \in \mathbf{W}\) 。设 \(a_1 \in \mathbf{K}^*\) 满足 \(v(a_1) = -\beta\) 。关系 \(ax_0 \in a_1^{-1} \mathbf{W}'\) 与 \(v(a) \leqslant \alpha\) 蕴含 \(a_1 ax_0 \in \mathbf{W}'\) 以及 \(v(a_0 a^{-1} a_1^{-1}) = \alpha + \beta - v(a) \geqslant \beta\) ，因而 \(a_0 x_0 = a_0 a^{-1} a_1^{-1}(a_1 ax_0) \in \mathbf{W}\) ，这是荒谬的；换言之，关系 \(ax_0 \in a_1^{-1} \mathbf{W}\) 蕴含 \(v(a) > \alpha\) 。

推论. 设 E 是 K 上的一个左拓扑向量空间，H 是 E 的一个闭超平面，D 是 E 的一个一维向量子空间且为 H 的代数余子空间。假设 v 不是非真的。那么 D 是 H 的拓扑余子空间。

考虑到命题 1 与命题 3，其证明与《拓扑向量空间》第一章 § 2 定理 1 的推论 2 相同。

命题 4. 假设 \(v\) 不是非真的且 \(K\) 是完备的。设 \(E\) 是 \(K\) 上的一个左拓扑向量空间，它是 Hausdorff 的且维数有限为 \(n\) 。对每个基 \((e_i)_{1 \leqslant i \leqslant n}\) （\(E\) 在 \(K\) 上的），映射 \((a_i) \mapsto \sum_{i=1}^{n} a_i e_i\) 由 \(K_s^n\) 到 \(E\) 上是拓扑向量空间同构。

考虑到命题 3 及其推论，其证明与《拓扑向量空间》第一章 § 2 定理 2 相同。

推论. 假设 v 不是非真的且 K 是完备的。设 E 是 K 上的一个 Hausdorff 拓扑向量空间，F 是 E 的一个有限维向量子空间。那么 F 是闭的。

F 是完备的。

